\cotefile{105}
\scanpage{19}
\handwritten{Nous \unsure{tr|to}}

% printed form
{\ttfamily\small
\begin{tabular}{@{}l@{\quad}l@{}}
/*JOB=PREDONTA,ACCT=*********,ALL=M06 & J0365CNU\\
//PREDONTA JOB MSGLEVEL=(1,1) & \\
 //\hspace{4em}EXEC CSSBATCH & 00000040\\
 //CRAYIN\hspace{1em}DD * & 00000050\\
 // & 00000160\\
 /* &
\end{tabular}}
% archive sheet number: 18

\scanpage{20}
\mspage{8}

We still have \deleted{to} see if $W\subset\widetilde W$. For the opposite inclusion we had to make restrictions on the $I$'s \inserted{(in \unsure{terms|favour} of $W$)} (condition 2), now to get $\widetilde W$ large enough, i.e. $TC$ large enough, we must make an assumption which will ensure that the bunch of $I$'s is large enough --- a natural condition would be that it contain a separating interval. We know already there \deleted{\unsure{are}}\inserted{\unsure{is one|are}} ``eligible'' \uline{separating} intervals, namely $W$-aspheric over $e$ and separating --- such as $L$.

% slanted
\margin{but this \unsure{cannot|can no} possibly be enough, \deleted{\unsure{it}} as it would (being fixed by itself) \deleted{\unsure{it}} give just one possible $\unsure{F'|F}$, hence one particular $\widetilde W$ being forced upon $W$....}

Let $f\in W$, consider factorization
\diagram{\begin{tikzcd}[ampersand replacement=\&,column sep=large]
X \arrow[r,"{\deleted{TC}\;i}"] \& X' \arrow[r,dash,"p"] \& Y \& i\in TC,\quad p\in F\quad\text{(or }i\in C,\;p\in TF\text{)}
\end{tikzcd}}
we know already that $i\in W$ (resp.! $p\in W$) \unsure{therefrom|therefore} $p\in W$ (resp. $i\in W$) --- hence we are reduced to \deleted{\unsure{the case $f$ is|Then one of}} proving either one of the two relations
\[
\begin{aligned}
 F\cap W&\subset TF &&\text{(necessary for }W\subset\widetilde W,\text{ as}\\
 &&& TF=F\cap\widetilde W\text{)}\\
 C\cap W&\subset TC &&(\mathrm{id})
\end{aligned}
\]
which are \inserted{\deleted{\unsure{both|not}}} equivalent \deleted{\unsure{each is|hence}} each to the condition $W\subset\widetilde W$ (or equivalently, $\deleted{\unsure{W=}}\widetilde W=W$). The second \inserted{[first?]} means that if $f:X-Y$ \inserted{[is in $W$ and]} has the \deleted{right} RLP with respect to $(TC)_0$ i.e. \inserted{to \deleted{\unsure{all}}} all $\alpha(\deleted{\unsure{i}}\unsure{i_1},i_0,I,e)$, then it has the RLP even with respect to \uline{all} monomorphisms which are $\in W$. \inserted{[????]}

We would like this to hold at least in the following two cases:
% slanted
\margin{In both cases $\mathscr A=A^{\wedge}$.}

a) $W=\underline W_A$ ($\underline W$ a basic localizer), $A$ a (local?) $\underline W$-test category

b) $A=A_0\times T$, \inserted{$\widehat A\simeq\underline{\mathrm{Hom}}(A_0^{\mathrm{opp}},T^\wedge)\simeq(A_0\unsure{\times}T)^\wedge$} with $T$ a test category, $W=\underline W_{A/A_0}$ ($\underline W$ a basic localizer, $\underline W_{A/A_0}$ is defined componentwise via $\underline W_T$.
% archive sheet number: 19

\scanpage{21}
% printed form
{\fontsize{5}{6}\selectfont
\begin{verbatim}
//*   PREDONTA0365        I
                           00000000      00000000      00000000      00000000
                         0000000000    0000000000    0000000000    0000000000
               00      0000  00      0000  00      0000  00      0000
              00     00 00  00     00 00  00     00 00  00     00 00
             00    00  00  00    00  00  00    00  00  00    00  00
            00   00   00  00   00   00  00   00   00  00   00   00
           00  00    00  00  00    00  00  00    00  00  00    00
          00 00     00  00 00     00  00 00     00  00 00     00
         0000      00  0000      00  0000      00  0000      00
        000       00  000       00  000       00  000       00
0000000000    0000000000    0000000000    0000000000
00000000      00000000      00000000      00000000

           PPPPPPPPPPP   RRRRRRRRRRR   EEEEEEEEEEEE  DDDDDDDDD     OOOOOOOOOOOO  NN        NN  TTTTTTTTTTTT   AAAAAAAAAA
           PPPPPPPPPPPP  RRRRRRRRRRRR  EEEEEEEEEEEE  DDDDDDDDDD    OOOOOOOOOOOO  NNN       NN  TTTTTTTTTTTT  AAAAAAAAAAAA
           PP        PP  RR        RR  EE            DD       DD   OO        OO  NNNN      NN       TT       AA        AA
                      PP        PP  RR        RR  EE            DD        DD  OO        OO  NN  NN    NN       TT       AA        AA
PP        PP  RR        RR  EE            DD        DD  OO        OO  NN NN     NN       TT       AA        AA
           PPPPPPPPPPPP  RRRRRRRRRRRR  EEEEEEEE      DD        DD  OO        OO  NN   NN   NN       TT       AAAAAAAAAAAA
           PPPPPPPPPPP   RRRRRRRRRRR   EEEEEEEE      DD        DD  OO        OO  NN    NN  NN       TT       AAAAAAAAAAAA
           PP            RR    RR      EE            DD        DD  OO        OO  NN     NN NN       TT       AA        AA
           PP            RR     RR     EE            DD        DD  OO        OO  NN      NNNN       TT       AA        AA
           PP            RR      RR    EE            DD       DD   OO        OO  NN       NNN       TT       AA        AA
           PP            RR       RR   EEEEEEEEEEEE  DDDDDDDDDD    OOOOOOOOOOOO  NN        NN       TT       AA        AA
           PP            RR        RR  EEEEEEEEEEEE  DDDDDDDDD     OOOOOOOOOOOO  NN         N       TT       AA        AA

****A  START  JOB  365  PREDONTA                        ROOM        9.20.10 AM 09 JUN 82  R6.PR1    SYS JES2  JOB  365  START  A****
****A  START  JOB  365  PREDONTA                        ROOM        9.20.10 AM 09 JUN 82  R6.PR1    SYS JES2  JOB  365  START  A****
****A  START  JOB  365  PREDONTA                        ROOM        9.20.10 AM 09 JUN 82  R6.PR1    SYS JES2  JOB  365  START  A****
****A  START  JOB  365  PREDONTA                        ROOM        9.20.10 AM 09 JUN 82  R6.PR1    SYS JES2  JOB  365  START  A****
****A  START  JOB  365  PREDONTA                        ROOM        9.20.10 AM 09 JUN 82  R6.PR1    SYS JES2  JOB  365  START  A****
****A  START  JOB  365  PREDONTA                        ROOM        9.20.10 AM 09 JUN 82  R6.PR1    SYS JES2  JOB  365  START  A****

DERNIERE MISE A JOUR DE JOURNAL: 08-06-82
STAGE FORMATION NLT-CLOTILDE LES 22 ET 23 JUIN
\end{verbatim}
}
% archive sheet number: 20

\scanpage{22}
\mspage{9}

Alternative presentation of theory, with possibly slightly less strong result --- as presentation of handy ns conditions on a $W$, so that there should exist a closed model structure on $\mathscr A$, s.th. $C=\mathrm{mono}$, $W=$ the given one, hence $TC=C\cap W$, $F$ defined by \deleted{lift.} RLP with respect to $TC=C\cap W$ \inserted{(say $(TC)_*$)}, $TF$ (``\unsure{Kan-injective|triv-injective}'' maps) defined by RLP with respect to $C$ \inserted{($TF=C_*$)}.

The problems which remain: (as we know \inserted{already} $(TF)^*=(C_*)^*=C$)

\textcircled{a} $\deleted{\unsure{Any|A_{\text{any}}}}F^*=((TC)_*)^*=TC$ --- and as $\deleted{\unsure{((TC)_*}}F^*\subset(TF)^*=C$, this amounts to: any monomorphism which has the LLP with respect to $F$ is in $W$. (A saturation type property of $W$ with respect to monomorphisms)

\textcircled{b} $TF=W\cap F$ --- in terms of the assumptions already made previously upon $W$, we know $TF\subset W$ i.e. $TF\subset W\cap F$, thus the problem is $W\cap F\subset TF$, i.e. any $f\in W$ which has the RLP with respect to $W\cap C$ is in $TF$ i.e. is Kan-injective.

c) Factorization $f=pi$ with either $p\in TF$, $i\in C$ (but this is already known and independent of $W$), or $p\in F$, $i\in TC$ \inserted{and even $i\in C\cap W$} --- and this should'nt be too hard in terms of 3) and 4) by standard arguments. Thus I'll admit c for the time being.

\textcircled{d} Moreover, I'm interested in the following, quite important for me, dual to 3): if $f\in F:X-Y$, and $(Y'-Y)\in W$, then $(X'-X)\in W$.

For handling \textcircled{a} \textcircled{b} \textcircled{d} presumably I'll have to rely strongly upon factorization c). Thus for a), using factorization of $f=pi$, and using $f\in F^*$ we get a section $s$ of $p$ \deleted{\unsure{f}} s.th. $sf=i$, thus $f$ is a retract of $i\in W$, and one wins if we assume 5).
\diagram{\begin{tikzcd}[ampersand replacement=\&,row sep=small,column sep=large]
\& X' \arrow[d,"{p\in F}"']\\
X \arrow[ur,hookrightarrow,"{i\in W\cap C}"] \arrow[r,hookrightarrow,"f"'] \& Y \arrow[u,bend right=75,"s"']
\end{tikzcd}}
% archive sheet number: 21

\scanpage{23}
% printed form
{\fontsize{5}{6}\selectfont
\begin{verbatim}
//*   PREDONT 0002   - CRAYI

                  PPPPPPPPPPP   RRRRRRRRRRR   EEEEEEEEEEEE  DDDDDDDDD     OOOOOOOOOOOO  NN        NN  TTTTTTTTTTTT
                  PPPPPPPPPPPP  RRRRRRRRRRRR  EEEEEEEEEEEE  DDDDDDDDDD    OOOOOOOOOOOO  NNN       NN  TTTTTTTTTTTT
                  PP        PP  RR        RR  EE            DD       DD   OO        OO  NNNN      NN       TT
                  PP        PP  RR        RR  EE            DD        DD  OO        OO  NN NN     NN       TT
                  PP        PP  RR        RR  EE            DD        DD  OO        OO  NN  NN    NN       TT
                  PPPPPPPPPPPP  RRRRRRRRRRRR  EEEEEEEE      DD        DD  OO        OO  NN   NN   NN       TT
                  PPPPPPPPPPP   RRRRRRRRRRR   EEEEEEEE      DD        DD  OO        OO  NN    NN  NN       TT
                  PP            RR    RR      EE            DD        DD  OO        OO  NN     NN NN       TT
                  PP            RR     RR     EE            DD        DD  OO        OO  NN      NNNN       TT
                  PP            RR      RR    EE            DD       DD   OO        OO  NN       NNN       TT
                  PP            RR       RR   EEEEEEEEEEEE  DDDDDDDDDD    OOOOOOOOOOOO  NN        NN       TT
                  PP            RR        RR  EEEEEEEEEEEE  DDDDDDDDD     OOOOOOOOOOOO  NN         N       TT

****A   END   JOB    2  PREDONT     - CRAY OUTPUT -     ROOM        9.20.01 AM 09 JUN 82  R6.PR1    SYS JES2  JOB    2   END   A****
****A   END   JOB    2  PREDONT     - CRAY OUTPUT -     ROOM        9.20.01 AM 09 JUN 82  R6.PR1    SYS JES2  JOB    2   END   A****
****A   END   JOB    2  PREDONT     - CRAY OUTPUT -     ROOM        9.20.01 AM 09 JUN 82  R6.PR1    SYS JES2  JOB    2   END   A****
****A   END   JOB    2  PREDONT     - CRAY OUTPUT -     ROOM        9.20.01 AM 09 JUN 82  R6.PR1    SYS JES2  JOB    2   END   A****
****A   END   JOB    2  PREDONT     - CRAY OUTPUT -     ROOM        9.20.01 AM 09 JUN 82  R6.PR1    SYS JES2  JOB    2   END   A****
****A   END   JOB    2  PREDONT     - CRAY OUTPUT -     ROOM        9.20.01 AM 09 JUN 82  R6.PR1    SYS JES2  JOB    2   END   A****
\end{verbatim}
}
% archive sheet number: 22

\scanpage{25}
\mspage{10}

$\mathscr A$ category

if $\Phi\subset\mathrm{Fl}(\mathscr A)$, let
\[
\begin{aligned}
\Phi^*&=\{u\in\mathrm{Fl}(\mathscr A)\mid u\text{ has the LLP with respect to }\Phi\}\\
\Phi_*&=\{u\in\mathrm{Fl}(\mathscr A),\;\rule{3em}{.3pt}\;\mathrm{RLP}\;\rule{5em}{.3pt}\;\}
\end{aligned}
\]

$\Phi$ is called left saturated if $\Phi=\Psi^*$ for some $\Psi$, i.e. $\Phi=(\Phi_*)^*$, \inserted{$\Psi$} right saturated if \deleted{it is} $\Psi=\Phi_*$ for some $\Phi$ i.e. $\Psi=(\Psi^*)_*$. Left saturation, right saturation $\Phi\subset\deleted{\unsure{\Phi}}(\Phi_*)^*$, $\Psi\subset(\Psi^*)_*$, associations $\Phi\mapsto\deleted{\unsure{\Psi^*}}\Phi_*$, $\Psi\mapsto\deleted{\unsure{\Phi_*}}\Psi^*$ are order reversing bijections between left saturated and right saturated subsets.
% slanted
\margin{Right saturation of $\Psi$ $\Rightarrow$ stability by composition, base change, retracts; dually for left saturation.}

The \uline{factorization property} for a pair $(\Phi,\Psi)$ of mutually orthogonal subsets of $\mathrm{Fl}(\mathscr A)$ ($\Phi\subset\Psi^*$ i.e. $\Psi\subset\Phi_*$), is the property \deleted{that} $\downarrow$
\[
\forall f\in\mathrm{Fl}(\mathscr A)\qquad \exists\ \text{factorization }f=pi,\quad\text{with }i\in\Phi,\ p\in\Psi
\]

\deleted{Let $\unsure{\overline\Phi,\overline\Psi}$ be the left and right saturations respectively}

\deleted{Prop.}\deleted{Let} \uline{Proposition} Assume factorization property holds \inserted{(for orthogonal pair $(\Phi,\Psi)$)}. Then $\Phi=\Psi^*$ iff $\Phi$ stable by direct factors. Dually $\Psi=\Phi_*$ iff $\Psi$ stable by retractions. Hence $(\Phi,\Psi)$ \deleted{are} dual pair of subsets of $\mathrm{Fl}(\mathscr A)$ (localizers) iff \deleted{and} \deleted{(they are orthogonal} both $\Phi$ and $\Psi$ are stable by direct factors.

\uline{Quillen pair} $(\Phi,\Psi)$:
\[
\begin{aligned}
 &\text{dual pair + factorization}\\
\Updownarrow\quad&\text{orthogonal pair + factorization + stability of }\Phi,\Psi\\
 &\hspace{12em}\text{by direct factors.}
\end{aligned}
\]

\deleted{Now let} \uline{Quillen triple} $(C,F,W)$ with a) $W$ \inserted{\deleted{\unsure{weakly}} mildly} saturated \textcircled{M5} (iso \deleted{\unsure{and}} in $W$, \deleted{and} \unsure{strong} compos., \uline{and stability by direct factors}.)

\margin{it should not be necessary here to assume the direct factor condition on $W$, if we assume b)}

b) $(C,F\cap W\overset{\mathrm{def}}=TF)$ and $(TC\overset{\mathrm{def}}=C\cap W,F)$ are Quillen pairs i.e. \textcircled{M1} \textcircled{M2} and \uline{stability of} \textcircled{$C$},\textcircled{$F$}, \deleted{\unsure{TC}} \inserted{$TC,TF$} \deleted{under direct factors} under direct factors (it is enough for $C,F$, provided true holds for \textcircled{$W$}).

\diagram{\(\begin{array}{cccccccc}
C&TF&&&&&&\\
\cup&\cap&W&\cdots&\cdots&\cdots&\cdots&\cdots\\
TC&F&&&&&&
\end{array}\)}

(hence \deleted{\unsure{$f$}} $f\in W\Longleftrightarrow f=pi,\ i\in TC,\ f\in TF$).

! This implies $W$ is strongly saturated i.e. if $\gamma_W:\mathscr A\to W^{-1}\mathscr A$, $f\in W\Longleftrightarrow \gamma_W(f)$ iso.
% slanted
\margin{true at any rate if in $\mathscr A$ finite direct and inverse limits exist (M0)}
% archive sheet number: 24

\scanpage{26}
% printed form
{\fontsize{5}{6}\selectfont
\begin{verbatim}
09:14:21      0.0000    CSP
09:14:21      0.0000    CSP      **********************************************************************
09:14:21      0.0000    CSP      *                                                                    *
09:14:21      0.0000    CSP      *                                                                    *
09:14:21      0.0000    CSP      *                                                                    *
09:14:21      0.0000    CSP      *         GETIA LE 07 JUIN 1982                                      *
09:14:21      0.0000    CSP      *         ---------------------                                      *
09:14:21      0.0000    CSP      *  ATTENTION                                                         *
09:14:21      0.0001    CSP      *  =========                                                         *
09:14:21      0.0001    CSP      *  LES FICHIERS NON ACCEDES DEPUIS 6 MOIS SERONT PURGES LE           *
09:14:21      0.0001    CSP      *  VENDREDI 04 JUIN 1982. CECI POUR LA CISI UNIQUEMENT.              *
09:14:21      0.0001    CSP      *                                                                    *
09:14:21      0.0001    CSP      *====================================================================*
09:14:21      0.0001    CSP      *  RAPPEL                                                            *
09:14:21      0.0001    CSP      *  ------                                                            *
09:14:21      0.0001    CSP      *  LES OPTIONS ON=0 ET ON=F DU COMPILATEUR FORTRAN CFT SONT TRES     *
09:14:21      0.0001    CSP      *  COUTEUSES EN TEMPS CPU, L'OPTION ON=0 EMPECHE TOUTE VECTORISATION-*
09:14:21      0.0002    CSP      *  ELLES NE DOIVENT DONC UTILISEES QU'A BON ESCIENT.                 *
09:14:21      0.0002    CSP      *                                                                    *
09:14:21      0.0002    CSP      *--------------------------------------------------------------------*
09:14:21      0.0002    CSP      *  RAPPEL    POUR LES UTILISATEURS EDF SEULEMENT                     *
09:14:21      0.0002    CSP      *  ------                                                            *
09:14:21      0.0002    CSP      *    LES PARAMETRES DE LA CARTE JOB DE L'ENVELOPPE IBM CSSBATCH SONT *
09:14:21      0.0002    CSP      *    DE 50K ET 1 SECONDE (OBLIGATOIRES)                              *
09:14:21      0.0002    CSP      *                 ------------------------------------               *
09:14:21      0.0002    CSP      *                                                                    *
09:14:21      0.0002    CSP      *     -RAPPEL DES NUMEROS DE TELEPHONE DU SUPPORT CRAY  :            *
09:14:21      0.0003    CSP      *       . EDF CLAMART  :  765-37-38                                  *
09:14:21      0.0003    CSP      *       . CISI SACLAY  :  908-64-87                                  *
09:14:21      0.0003    CSP      *       . CISI RUEIL   :  732-02-73 POSTE 552                        *
09:14:21      0.0003    CSP      *                                                                    *
09:14:21      0.0003    CSP      *                                                                    *
09:14:22      0.0003    CSP      **********************************************************************
09:14:22      0.0003    CSP
09:14:22      0.0003    CSP
09:14:22      0.0003    CSP             CRAY-1S/1000   SERIAL-26        GETIA-CLAMART           06/09/82
09:14:22      0.0003    CSP
09:14:22      0.0004    CSP             CRAY-1  OPERATING SYSTEM        COS 1.10  ASSEMBLY DATE 05/11/82
09:14:22      0.0004    CSP
09:14:22      0.0004    CSP
09:14:22      0.0005    CSP     JOB,JN=PREDONTA,T=9,TIO=9.
09:14:22      0.0009    CSP     ACCOUNT,AC=*********.
09:14:22      0.0012    CSP     CFT.
09:14:22      0.0015    USER    CF000 - CFT VERSION -   05/10/82 1.09
09:14:22      0.0035    USER    CF001 - COMPILE TIME =   0.0020 SECONDS
09:14:22      0.0035    USER    CF002 -        5 LINES,        5 STATEMENTS
09:14:22      0.0036    CSP     LDR.
09:14:23      0.0235    USER    LD000 - BEGIN EXECUTION
09:14:23      0.0236    USER    FT063 -    STOP              IN GUTS
09:14:23      0.0236    CSP     END OF JOB
09:14:23      0.0236    CSP
09:14:23      0.0236    CSP
09:14:23      0.0236    CSP
09:14:23      0.0236    CSP             JOBNAME PREDONT SQ 30624
09:14:23      0.0236    CSP             TIME EXECUTING IN CPU - 00:00:00.0236
09:14:23      0.0237    CSP             TIME WAITING TO EXECUTE 00:00:00.7018
09:14:23      0.0237    CSP             TIME WAITING FOR I/O -- 00:00:00.7687
09:14:23      0.0237    CSP             MEMORY USAGE ---     0.02974 MWDS-SEC
09:14:23      0.0237    CSP             DISK BLOCKS MOVED -----      259
09:14:23      0.0237    CSP             PHYSICAL I/O REQUESTS -       32
\end{verbatim}
}
% archive sheet number: 25

\scanpage{27}
$Y$ is Kan

\diagram{\begin{tikzcd}[ampersand replacement=\&,row sep=small,column sep=large]
X\arrow[r,"\text{fibration}"'] \& Y\arrow[dr,dash,bend left=12]\\
\&\&\text{Kan}
\end{tikzcd}}

\deleted{$\underline{\mathrm{Hom}}(I,Y)\longrightarrow Y\times Y$ is a Kan fibration (if $Y$ is Kan)

is it a Serre fibration

\[
a\longrightarrow Y\times Y\qquad \alpha,\beta\in Y_a\qquad
\underline{\mathrm{Hom}}_{\alpha,\beta}(I_a,Y_a)\qquad
\underline{\mathrm{Hom}}_{\alpha_b,\beta_b}(I_b,Y_b)
\]

a) $\mathrm{Iso}_M\subset W$

b) If $f=\unsure{vu|uv}$ and two among $u,v,f$ are in $W$, so is the third $\big)\Rightarrow$ \inserted{mild} saturation

c) $W$ stable by retraction \deleted{(\unsure{other question|stronger|stronger than} $\longrightarrow$}

\margin{but does it imply saturation namely $fg\in W,\ gf\in W\Longrightarrow f,g\in W$??}

\diagram{\begin{tikzcd}[ampersand replacement=\&,row sep=small,column sep=large]
A\arrow[r,shift left,"f"] \& B\arrow[l,shift left,"g"] \& fg\in W,\quad gf\in W\Longrightarrow f,g\in W\;?
\end{tikzcd}}

\[
fg\in\mathrm{id}_A,\quad gf\unsure{\in|=|\leftrightarrow}\mathrm{id}_A\Longrightarrow f,g\in W
\]

\diagram{\begin{tikzcd}[ampersand replacement=\&,row sep=small,column sep=large]
A\arrow[r,"{\deleted{\unsure{fg}}gf}"] \& A \& A\arrow[r,"{\deleted{\unsure{fg}}=}"] \& A\\
A\arrow[r,"f","{\unsure{\sigma}}"'] \& B\\
B\arrow[r,"fg"] \& B\\
A \& B\arrow[l,"f"']\arrow[l,bend right=30,"{\unsure{g}}"]
\end{tikzcd}}

$\deleted{\unsure{gf}}fg$

\diagram{\begin{tikzcd}[ampersand replacement=\&,row sep=small,column sep=large]
A\arrow[r,"f"]\arrow[d,"f",shift right] \& B\arrow[d,"{=}",shift left]\\
B\arrow[u,"g",bend left=55]\arrow[u,shift right]\arrow[r,"fg","{\in W}"'] \& B\arrow[u,shift left]
\end{tikzcd}}

\[
\boxed{\begin{array}{cccc}a)&b)&c_1)&\deleted{\unsure{d}}\\ &&c_2)&\\&&c_3)&\end{array}}
\]

\diagram{\begin{tikzcd}[ampersand replacement=\&,row sep=small,column sep=large]
A\arrow[r,dash,"f"]\arrow[d,dash,"i"'] \& B\arrow[d,dash,"j"]\\
 A'\arrow[u,"p",bend left=25]\arrow[r,dash,"f'"] \& B'\arrow[u,"q",bend left=25]
\end{tikzcd}}

\[
\left|\begin{aligned}f'i&=jf\\fp&=qf'\\pi&=\mathrm{id}_A\\qj&=\mathrm{id}_B\\f'&\in W\end{aligned}\right.\qquad\overset{?}{\Longrightarrow}\quad f\in W.
\]

\diagram{\begin{tikzcd}[ampersand replacement=\&,row sep=small,column sep=large]
A\arrow[d,"f"] \& W \& \unsure{h_W|hW|{}^hW}\\ B\arrow[u,bend left=20,"g"]
\end{tikzcd}}

\[
\begin{aligned}
gf&\underset h\sim\mathrm{id}_A&\text{donc }gf&\in W\\
fg&\underset h\sim\mathrm{id}_B&fg&\in W
\end{aligned}
\]

\diagram{\begin{tikzcd}[ampersand replacement=\&,row sep=small,column sep=large]
\& \unsure{L|\lambda}\arrow[dl,dash,"{\unsure{d}}"']\arrow[dr,dash] \& \& L\arrow[r,bend left,"{\unsure{7|?}}"]\&X\arrow[l,"d"]\\
 X\&\& I\times X:y\arrow[ll,"{\unsure{sg}}"']\arrow[ll,bend right=20,"{\unsure{o}}"]
\end{tikzcd}}

\diagram{\begin{tikzcd}[ampersand replacement=\&,row sep=small,column sep=large]
\&\&L\\
\& X\arrow[r,"{\unsure{M}}","{\unsure{o}}"']\&X\arrow[u,"d"]\arrow[l,bend left=40]\arrow[l,bend left=65]\\
X\times X\arrow[ur,"{\unsure{y}\deleted{\unsure{bo}}}"]\&\deleted{\unsure{g}}\&
\end{tikzcd}}

(i) $gf=\mathrm{id}_A,\ fg\in W\Longrightarrow f,g\in W$

(ii) $gf\in W,\ fg\in W\Longrightarrow f,g\in W$

(iii) $W$ stable \deleted{by \unsure{facteurs|factors}} direct factors

\diagram{\begin{tikzcd}[ampersand replacement=\&,row sep=small,column sep=large]
\text{(ii)}\arrow[dr,Rightarrow]\&\&\text{(iii)}\arrow[dl,Rightarrow]\\
\&\text{(i)}
\end{tikzcd}}}
% archive sheet number: 26

\scanpage{28}
\mspage{10 bis}

\handwritten{\[
\begin{aligned}
X\xrightarrow{f}Y\qquad &X,Y\in M_{c,f}\\ &f\in F
\end{aligned}
\]
then $\underset{\text{i.e.}}{f\in TF}\Longleftrightarrow\text{retract property}\Longleftrightarrow$

\diagram{\(\begin{array}{cc}C&F\cap W=TF\\ \cup&\cap\\ TC=C\cap W&F\end{array}\)}

\[
\text{\textcircled{a}}\left\{\begin{aligned}C&=(TF)^*\\F&=(TC)_*\end{aligned}\right.
\qquad \text{\textcircled{b}}\quad
\deleted{\unsure{W'}}W=\{pi\mid \deleted{\unsure{p\in TF}},\
p\in\underset{\widetilde{TF}}{\underset{\Vert}{C_*}},\
i\in\underset{\overline{TC}}{\underset{\Vert}{F^*}}\deleted{\unsure{,i}}\}
\]

$\swarrow$

$\widetilde{TF}\subset F,\ \overline{TC}\subset C$ (because $F,C$ \inserted{\unsure{r. and l.|are l.|trivially}} saturated)

\deleted{$\left.\begin{aligned}TF&=\widetilde{TF}\\ TC&=\overline{TC}\end{aligned}\right\}$ because $TF=F\cap W,\ TC=C\cap W$ are stable under direct factors. [provided we prove $W$ is stable]}

we only have to \deleted{prove}\inserted{express}, to get b in terms of a), that \deleted{\unsure{if}}
\[
\boxed{\widetilde{TF}\subset W,\ \overline{TC}\subset W}
\]
and this is equivalent with
\[
\boxed{\widetilde{TF}=TF,\ \overline{TC}=TC}
\]
i.e. $TF,TC$ saturated i.e. $(C,TF),(TC,F)$ Quillen pairs

\uline{Equivalent conditions in a model category} (in Quillen's sense)

\deleted{Closed model}

(i) satisfies \textcircled{a} and \textcircled{b}, i.e closed model category in Quillen sense

$\Updownarrow$ (ii) $(C,TF=F\cap W)$ and $(TC=C\cap W,F)$ are Quillen pairs

$\Uparrow$ (iii) $C,F,W$ are stable under \unsure{direct|retract} factors.

But I've to check that (i)(ii) (which are equivalent) does imply that \uline{$W$} is stable under direct factors. But Quillen proves it even implies $W$ is strongly saturated (i.e $f\in W$ iff $\gamma(f)$ iso) --- at least under the assumption that finite direct and inverse limits exist in $M$.

for any model category

\diagram{\begin{tikzcd}[ampersand replacement=\&,row sep=small,column sep=large]
\&\&X_2\arrow[r,"{C\text{ or }F}","g"']\arrow[drr,dashed,bend left=9,"F"] \&Y_2\arrow[dr,"TF"]\\
\varnothing\arrow[r,"C"']\&X_1\arrow[ur,"TC"]\arrow[urr,dashed,bend right=8,"C"']\arrow[dr,"TF"']\&\&\&Y_1\arrow[r,"F"']\& e\\
\&\&X\arrow[r,"f"]\&Y\arrow[ur,"TC"']
\end{tikzcd}}

\[
\begin{aligned}\gamma(f)\ \mathrm{iso}&\Longleftrightarrow\gamma(g)\ \mathrm{iso}\\
f\in W&\Longleftrightarrow g\in W\end{aligned}
\quad\left|\quad
\begin{aligned}
&W\text{ strongly saturated}\\
\Longleftrightarrow{}&\text{for }f:X\to Y,\;X,Y\in M_{cf},\\
&f\in C,\ \gamma(f)\mathrm{iso}\Rightarrow f\in W.\\
\Longleftrightarrow{}&\text{for }f:X-Y,\;X,Y\in M_{cf},\\
&f\in F,\ \gamma(f)\mathrm{iso}\Rightarrow f\in W.
\end{aligned}\right.
\]}

% printed form
\mspage{1}
{\fontsize{5}{6}\selectfont
\begin{verbatim}
      PAGE 1                ON=CELMPQRSTV              06/09/82-09:14:22      CFT 1.09(05/10/82)  PAGE 1
1     1.        PROGRAM GUTS
2     2.        WRITE(6,1)
3     3.      1 FORMAT(' BONJOUR')
4     4.        STOP
5     5.        END
\end{verbatim}
}
% archive sheet number: 27

\scanpage{30}
\mspage{11}

To construct any Quillen triple, we may proceed as follows: Start with \inserted{a localizer} $C\subset\mathrm{Fl}(\mathscr A)$, \textcircled{a} \deleted{\unsure{being}}\inserted{$C$ stable} \deleted{\unsure{$C$}} by direct factors, hence a pair $(C,C_*=TF)$ --- we \uline{must check \textcircled{b} factorization} \uline{condition for the pair} $(C,C_*)$. \deleted{\unsure{Now|Hence}} Thus $(C,TF)$ is a Quillen pair.

Let next \textcircled{c} $\boxed{W\supset TF}$ any localizer which is \inserted{\textcircled{d} \uline{$W$ mildly}} \uline{saturated and stable under direct factor}. We now get $TC=C\cap W$ \inserted{$\subset C$}, which is \deleted{\unsure{again}} stable under direct factor. Take $F=(TC)_*$ \inserted{$(\supset TF=C_*)$}, we must check \textcircled{\unsure{e|d}} \uline{factorization property for the pair} $(TC,F)$ --- thus we get that $(TC,F)$ is a Quillen pair. Lastly, we must still check that $TF=F\cap W$, and by c) we know already $TF\subset F\cap W$, we have to check still that $F\cap W\subset TF$. But using factorization $f=pi,\ p\in TF,\ i\in C$, as $f,p\in W$ we get $i\in W$ hence $i\in W\cap C=TC$, \deleted{hence} there exists $s:X'-X$ s that $si=\mathrm{id}_X$ (retraction of $X'$ upon $X$) and $fs=p$, hence $f$ a retraction of $p$, hence is in $TF$, OK. Thus

\diagram{\begin{tikzcd}[ampersand replacement=\&,row sep=small,column sep=large]
X\arrow[r,dash,"="]\arrow[d,"{i\in TC}"'] \& X\arrow[d,"{f\in F}"] \& \& X'\arrow[r,shift left,"s"]\arrow[d,dash,"p"'] \& X\arrow[l,bend left=20,"i"]\arrow[d,dash,"f"]\\
X'\arrow[r,"p"']\arrow[ur,dashed,"s"]\&Y \&\&Y\arrow[r,shift left,"="]\&Y\arrow[l,shift left]
\end{tikzcd}}

\uline{Proposition} In order for a pair \inserted{of localizers} $(C,W)$ to correspond to a Quillen triple $(C,F,W)$, it is ns that the following conditions be satisfied
\[
\left\{\begin{aligned}
1)&\ C,W\ \text{\uline{stable by direct factors}},\ \deleted{\unsure{+}}\ W\text{ mildly saturated}\\
2)&\ \text{Factorization condition for }(C,C_*)\text{ and for }(C\cap W,(C\cap W)_*).\\
3)&\ W\supset TF\overset{\mathrm{def}}=C_*
\end{aligned}\right.
\]

We have to develop mainly a criterion to get factorization. [In the case of interest to us, we get 1) free, 3) almost free (there comes in the ``test''-assumption though), and 2) causes a little technical problem. I am however more than in factorization conditions, interested in the
% archive sheet number: 29

\scanpage{31}
\mspage{2}
% printed form
{\fontsize{5}{6}\selectfont
\begin{verbatim}
GUTS      PAGE 2                  ON=CELMPQRSTV                06/09/82-09:14:22      CFT 1.09(05/10/82)  PAGE 2       GUTS

TABLE OF STATEMENT NUMBERS (ALL ADDRESSES IN TABLES ARE IN OCTAL)
--------------------------

  NUMBER USE

      1 -FN

  (SN=STATEMENT NUMBER, GSN=GENERATED STATEMENT NUMBER)
  (FN=FORMAT NUMBER, UNDEF*=UNDEFINED STATEMENT NUMBER)

TABLE OF NAMES ENCOUNTERED (ADDRESS FOR DUMMY ARGUMENT IS THE ARGUMENT NUMBER)
--------------------------

 ADDRESS NAME    TYPE MAIN USAGE  BLOCK      ADDRESS NAME    TYPE MAIN USAGE  BLOCK      ADDRESS NAME    TYPE MAIN USAGE  BLOCK

         $MAIN       ENTRY                           $WFF        EXTERNAL                      12 GUTS     R  ENTRY
         $STOP       EXTERNAL                        $WFI        EXTERNAL

 BLOCK NAMES AND LENGTHS IN OCTAL

      23-GUTS              4-#TB
BONJOUR
\end{verbatim}
}
% archive sheet number: 30

\scanpage{32}
\mspage{12}
following properties (characterizing a \uline{strict} Quillen triple)

4) \deleted{A any cobase change} \inserted{$W$ stable by base} change by any $\deleted{\unsure{g}}f\in F$, and by cobase change by any $\deleted{\unsure{f}}g\in C$.

\deleted{\unsure{5)}} \uline{NB} From my point of view, $W$ is given beforehand in a rather tangible way, by \unsure{cohomological|homological} criteria, which turns out \unsure{very|easy} handy. In Kan's--Quillen's point of view, $W$ was a little more hidden maybe in terms of (say) ss structures, they get to it via \deleted{\unsure{a}} the pair $(C,TF)$ and $F\supset TF$ the Kan fibrations, which is a somewhat ad-hoc construction in terms of a certain \deleted{\unsure{g}} a-priori generating set $(TC)_0$ --- corresponding roughly to expressing the horn lifting property. What remains to be understood, in the context say of general test categories, is the relationship between this and the direct description $TC=C\cap W$, \deleted{\unsure{what}} i.e why this is just the \uline{saturation} of $(TC)_0$. Maybe I should have a look at Gabriel-Zisman and their treatment of ``anodyne extensions''.
% archive sheet number: 31

\scanpage{33}
% printed form
{\fontsize{5}{6}\selectfont
\begin{verbatim}
//*   PREDONT 0002  - CRAYI

                  PPPPPPPPPPP   RRRRRRRRRRR   EEEEEEEEEEEE  DDDDDDDDD     OOOOOOOOOOOO  NN        NN  TTTTTTTTTTTT
                  PPPPPPPPPPPP  RRRRRRRRRRRR  EEEEEEEEEEEE  DDDDDDDDDD    OOOOOOOOOOOO  NNN       NN  TTTTTTTTTTTT
                  PP        PP  RR        RR  EE            DD       DD   OO        OO  NNNN      NN       TT
                  PP        PP  RR        RR  EE            DD        DD  OO        OO  NN NN     NN       TT
                  PP        PP  RR        RR  EE            DD        DD  OO        OO  NN  NN    NN       TT
                  PPPPPPPPPPPP  RRRRRRRRRRRR  EEEEEEEE      DD        DD  OO        OO  NN   NN   NN       TT
                  PPPPPPPPPPP   RRRRRRRRRRR   EEEEEEEE      DD        DD  OO        OO  NN    NN  NN       TT
                  PP            RR    RR      EE            DD        DD  OO        OO  NN     NN NN       TT
                  PP            RR     RR     EE            DD        DD  OO        OO  NN      NNNN       TT
                  PP            RR      RR    EE            DD       DD   OO        OO  NN       NNN       TT
                  PP            RR       RR   EEEEEEEEEEEE  DDDDDDDDDD    OOOOOOOOOOOO  NN        NN       TT
                  PP            RR        RR  EEEEEEEEEEEE  DDDDDDDDD     OOOOOOOOOOOO  NN         N       TT

****A  START  JOB    2  PREDONT     - CRAY OUTPUT -     ROOM        9.19.44 AM 09 JUN 82  R6.PR1    SYS JES2  JOB    2  START  A****
****A  START  JOB    2  PREDONT     - CRAY OUTPUT -     ROOM        9.19.44 AM 09 JUN 82  R6.PR1    SYS JES2  JOB    2  START  A****
****A  START  JOB    2  PREDONT     - CRAY OUTPUT -     ROOM        9.19.44 AM 09 JUN 82  R6.PR1    SYS JES2  JOB    2  START  A****
****A  START  JOB    2  PREDONT     - CRAY OUTPUT -     ROOM        9.19.44 AM 09 JUN 82  R6.PR1    SYS JES2  JOB    2  START  A****
****A  START  JOB    2  PREDONT     - CRAY OUTPUT -     ROOM        9.19.44 AM 09 JUN 82  R6.PR1    SYS JES2  JOB    2  START  A****
****A  START  JOB    2  PREDONT     - CRAY OUTPUT -     ROOM        9.19.44 AM 09 JUN 82  R6.PR1    SYS JES2  JOB    2  START  A****

DERNIERE MISE A JOUR DE JOURNAL: 08-06-82
STAGE FORMATION NLT-CLOTILDE LES 22 ET 23 JUIN
\end{verbatim}
}
% archive sheet number: 32

\scanpage{34}
\mspage{13}
% slanted
\margin{\uline{Factorization} \unsure{sy.}\ \uline{Problem} of (cf th. p.\ 16, which gives the opposite result\ldots\ \deleted{\uline{NB} we need only the})}
$M$ a category, $\Phi_0\subset\mathrm{Fl}(M)$, $\Psi=(\Phi_0^{\deleted{\unsure{*}}})_*$, $\Phi_0\subset\Phi\subset\overline{\Phi}_0=\Psi^*$, \inserted{NB. we don't use $\Phi\in\unsure{\Phi_0^*|\overline{\Phi}_0}$} we want to deduce conditions on $\Phi_0,\Phi$ insuring that the pair $(\Phi,\Psi)$ satisfies factorization condition, i.e. $\forall f\in\mathrm{Fl}(M)$, $\exists$ factorization $f=pi$, $i\in\Phi$, $p\in\Psi$.

The idea is to take \inserted{(for given $f$)} all \inserted{comm.} diagrams
\diagram{\[
\tag*{(D)}
\begin{tikzcd}[ampersand replacement=\&]
A \arrow[r,"v"] \arrow[d,"g"'] \& X \arrow[d,"f"]\\
B \arrow[r,"u"'] \& Y
\end{tikzcd}
\]}
giving rise to
\diagram{\begin{tikzcd}[ampersand replacement=\&,row sep=small,column sep=large]
A \arrow[r,"v"] \arrow[d,"\alpha"',"\mathrm{coc}" near end] \& X \arrow[d,"\deleted{\unsure{g}}i_D"] \arrow[dd,bend left=56,"f"]\\
B \arrow[r] \& X_D=X\amalg_A B \arrow[d,"p_D"]\\
\&Y
\end{tikzcd}}
with $g\in\Phi_0$ and to define
\[
\deleted{\unsure{\Sigma_1(f)|Z_1(f)}}\inserted{X_1(f)}
=\deleted{\unsure{\lim|\varinjlim}}\coprod_{\text{all }D}{}_{(X)}X_D
\qquad(\deleted{\unsure{am}} \text{amalgamated sum under }X)
\]
so that we get factorization
\diagram{\begin{tikzcd}[ampersand replacement=\&,row sep=small,column sep=large]
X \arrow[r,"i_1(f)"] \arrow[d,"f"'] \& \deleted{\unsure{Z|\Sigma}}X_1(f) \arrow[dl,"\Sigma_1(f)"]\\
Y \&
\end{tikzcd}}
with (hopefully) $i_1(f)\in\Phi$.

For this first step, we used

a) $\Phi_0$ is small (we'll see later how to relax this condition)

b) Amalgamated sums indexed by $I$ with $\mathrm{card}\,I\leq\mathrm{card}\,\Phi_0$ \deleted{\unsure{set}} \inserted{rather, $\overbrace{\mathrm{card}\,\Phi_0\times c}^{\delta}$} exist in $M$, \deleted{and} \inserted{where $c$ is $\sup(\mathrm{card}\,\mathrm{Hom}(A,X)\times\deleted{\unsure{\mathrm{card}}}\mathrm{Hom}(B,Y))$ for all $u:A\to B$ in $\Phi_0$}

c) $\Phi$ stable by \inserted{composition, and by} cobase extension, and (by \inserted{so} filtering direct limits indexed by a set of cardinality $\leq\mathrm{card}\,\Phi_0$ \deleted{\unsure{etc}} \inserted{rather $\leq\unsure{d|\delta|c}$})
% slanted
\margin{\unsure{do|$d_0$} will be strengthened $\downarrow$}

We then are writing a transfinite induction, to define

$X_\alpha(f)$ and
\diagram{\begin{tikzcd}[ampersand replacement=\&,row sep=small,column sep=large]
X \arrow[r,"i_\alpha(f)"] \arrow[d,"f"'] \& X_\alpha(f) \arrow[dl,"\Sigma_\alpha(f)"]\\
Y\&
\end{tikzcd}}
$i_\alpha(f)\in\Phi$ with $\alpha$ any ordinal $\deleted{\unsure{\beta}}<\alpha_0$ ($\alpha_0$ \deleted{first} \deleted{of} ordinal (\deleted{\unsure{of cardinality strictly}} \deleted{\unsure{with cardinality at least}}) \inserted{\unsure{number|namely|suitable} to be fixed} \deleted{\unsure{fix}} later), \deleted{defined later in terms of $\kappa$},

\diagram{\[
\left\{\begin{aligned}
\Sigma_{\alpha+1}(f)&=\Sigma_1(\Sigma_\alpha(f)),\quad X_{\alpha+1}(f)=\text{source }\Sigma_{\alpha+1}(f),\quad i_{\alpha+1}(f)=\text{composition}\\
&\begin{tikzcd}[ampersand replacement=\&,column sep=large]
X \arrow[r,"i_\alpha(f)"] \arrow[d,"f"'] \& X_\alpha(f) \arrow[r,"i_1(\Sigma_\alpha(f))"] \arrow[dl,"\Sigma_\alpha(f)"'] \& X_{\alpha+1}(f)\arrow[dll,"\Sigma_{\alpha+1}(f)=\Sigma_1(\Sigma_\alpha(f))"]\\
Y\&\&
\end{tikzcd}\\
\Sigma_\alpha(f)&=\varinjlim_{\alpha'<\alpha}\Sigma_{\alpha'}(f)\qquad\text{if $\alpha$ is a limiting ordinal.}
\end{aligned}\right.
\]}

\deleted{\unsure{We|NB}} We assume that $\Phi$ is stable too by \inserted{the} filtering direct limits \deleted{of} \deleted{ordinality} which occur in the second transfinite induction step.

Let $\pi$ be a cardinal such that \deleted{all objects of $\Phi_0$} \inserted{for any $g\in\Phi_0$, source and target of $g$} are $\pi$-accessible --- i.e. $\mathrm{Hom}(A,-):M\to\mathrm{Sets}$ commutes to filtering direct limits which are \deleted{\unsure{greater w.|great wr}} ``large with resp to $\pi$''. We may take $\pi=\sup_{g\in\Phi_0}\pi_g$, where
% archive sheet number: 33

\scanpage{35}
% printed form
\diagram{{\footnotesize\ttfamily\setlength{\tabcolsep}{4pt}
\begin{tabular}{@{}llllllll@{}}
IMTADE & IMTAE & IMTASD & IMTA01 & IMTYNO & IMVISU & IMVVPR & IM1BDC \\
INICSD & INITMP & INTAB0 & ISNOP5 & LEAMAT & LEAMMS & LEATRI & LEB \\
LEBDCL & LECOOR & LEFORC & LEHYMA & LEMAIL & LEMEFI & LEMILI & LEMPLR \\
LEMUA & LENOPO & LETAE & LETASD & LETYNO & LEVISU & LEVVPR & LRNDL1 \\
LSVVPR & MONOPO & MONOP5 & MONOP6 & NCMUNS & NDL & NIVLIB & NUPAG \\
NUPAGE & PRINTB & PRIN1B & RACRNO & REAMAT & REATRI & REB & REBDCL \\
RECOOR & REFORC & REHYMA & REMAIL & REMEFI & REMILI & REMPLR & REMUA \\
RENOPO & RESOUR & RETAE & RETASD & RETYNO & REVISU & REVVPR & RSTSDE \\
SAUCSD & SAUSDS & SAUTP & SDBMSC & STAMD & STAMR & TRAMAD & TRAMAR \\
TUERSD & TUTASD & VERNN & VERTNN \\
\hline
IMAMA1 & IMAMA2 & \tikz[remember picture,overlay]\coordinate (codeb);IMATRI & IMB & IMBDCL & IMB1 & \tikz[remember picture,overlay]\coordinate (codef);IMCOOR & \unsure{IMCOO1|IMCO01} \\
IMFORC & IMHYMA & IMMAIL & IMMAI1 & \tikz[remember picture,overlay]\coordinate (coded);IMMILI & IMMPLR & IMMUA & IMNOP0\tikz[remember picture,overlay]\coordinate (codeh); \\
IMTADE & IMTAE & IMTASD & IMTA01 & IMTYNO & IMVISU & IMVVPR & IM1BDC \\
INICSD & INITMP & INTAB0 & ISNOP5 & LEAMAT & LEAMMS & LEATRI & LEB \\
LEBDCL & LECOOR & LEFORC & LEHYMA & LEMAIL & LEMEFI & LEMILI & LEMPLR\tikz[remember picture,overlay]\coordinate (codei); \\
LEMUA & LENOPO & LETAE & LETASD & LETYNO & LEVISU & LEVVPR & LRNDL1 \\
LSVVPR & MONOPO & MONOP5 & MONOP6 & NCMUNS & NDL & NIVLIB & NUPAG \\
NUPAGE & PRINTB & PRIN1B & RACRNO & REAMAT & REATRI & REB & REBDCL \\
RECOOR & REFORC & REHYMA & REMAIL & REMEFI & REMILI\tikz[remember picture,overlay]\coordinate (codeg); & REMPLR & REMUA \\
\tikz[remember picture,overlay]\coordinate (codea);RENOPO & RESOUR & RETAE & RETASD & RETYNO & REVISU & REVVPR & RSTSDE \\
SAUCSD & SAUSDS & \tikz[remember picture,overlay]\coordinate (codec);SAUTP & SDBMSC\tikz[remember picture,overlay]\coordinate (codee); & STAMD & STAMR & TRAMAD & TRAMAR \\
TUERSD & TUTASD & VERNN & VERTNN \\
\end{tabular}}}
\handwritten{\figure{\begin{tikzpicture}[remember picture,overlay]
\draw[thin] ([xshift=-2pt,yshift=-2pt]codea)
 .. controls +(5pt,-2pt) and +(-16pt,2pt) .. ([xshift=5pt,yshift=4pt]codeb)
 .. controls +(11pt,2pt) and +(-11pt,-4pt) .. ([yshift=2pt]codec)
 .. controls +(13pt,0pt) and +(-12pt,4pt) .. ([xshift=2pt,yshift=4pt]coded)
 .. controls +(10pt,1pt) and +(-10pt,-14pt) .. ([xshift=2pt,yshift=-2pt]codee)
 .. controls +(14pt,-4pt) and +(-15pt,-2pt) .. ([xshift=4pt,yshift=4pt]codef)
 .. controls +(13pt,4pt) and +(-11pt,-8pt) .. ([xshift=-5pt,yshift=-2pt]codeg)
 .. controls +(10pt,-3pt) and +(-6pt,6pt) .. (codeh)
 .. controls +(2pt,-4pt) and +(3pt,5pt) .. ([xshift=-1pt]codei);
\end{tikzpicture}}}
% archive sheet number: 34

\scanpage{36}
\mspage{14}
$\pi_g$ is the smallest cardinal s.th. source and target of $g$ are $\pi$-accessible.

\uline{Example} If the objects \unsure{occuring|occurring} as sources and targets \deleted{of} $g$'s in $\Phi_0$ are ``of finite presentation'', i.e. $\mathrm{Hom}(A,-)$ commutes with any filtering direct limit, we get $\pi=0$. In this case we stop the induction with $\alpha_0=$ first infinite ordinal \inserted{$\omega$} \deleted{the first}. In case the objects $A,B$ are not of finite presentation (\inserted{\unsure{we assume|we can take|we mean}} still they are accessible --- we take the smallest cardinal such that $\mathrm{Hom}(A,?)$ commutes with direct limits indexed by \deleted{ordi} ordered sets \deleted{great} \inserted{large} w.r.t.\ $\pi$. Let \inserted{then} $\alpha_0$ be the first ordinal s.th. the \unsure{interval|ordset} $I_{\alpha_0}=\{\text{all ordinals }\alpha\text{ with }\alpha<\alpha_0\}$ be \deleted{great} large wrt $\pi$. We take the factorization

\diagram{\begin{tikzcd}[ampersand replacement=\&,row sep=small,column sep=large]
X \arrow[r,"i_{\alpha_0}(f)=i(f)"] \arrow[d] \& X_{\alpha_0}(f)=\deleted{\unsure{Z|\Sigma}}X(f) \arrow[dl,"\Sigma_{\alpha_0}(f)=p(f)"]\\
Y\&
\end{tikzcd}}
and we get thus $i(f)\in\Phi$, $p(f)\in\Psi$, provided we make the following assumptions;

d) The \deleted{\unsure{elem|objects}} source and target objects of the arrows of $\Phi_0$ are \uline{accessible} (ok if $M$ a topos etc ---) and, if $\pi_0$ is the cardinality of the smallest ordinal number \inserted{$\alpha_0$} s.th. for any such $A$, $\mathrm{Hom}(A,-)$ commutes to the \deleted{\unsure{f}} direct limit of type $I_{\alpha_0}$, \deleted{\unsure{we|we assume}} and \deleted{$\pi=\sup(\pi$} $\pi=\sup(\pi_0,\overset{\inserted{\text{rather, }\delta}}{\overbrace{\mathrm{card}\,\Phi_0}})$, we assume that for any filtering direct system with indexing set $I$ of cardinality $\leq\pi$, and transition maps in $\Phi$, the direct limit exists and $\deleted{\unsure{X_i}} Z_i\to Z$ are equally in $\Phi$.

Thus we get, not only factorization $f=pi$ with $i\in\Phi$, $p\in\Psi$, but even factorization $\deleted{\unsure{\Sigma|Z}}X(f)$ \uline{functorial} with regard to variable $f$.
% archive sheet number: 35
